\documentclass[12pt]{article}
\usepackage[utf8]{inputenc}
\usepackage[T1]{fontenc}
\usepackage{hyperref}
\title{Wagiman Language Resources for K-12 Curriculum Development}
\author{Compiled by Claude Sonnet 5 (claude-sonnet-5) for Andr\'as Kornai}
\date{}
\begin{document}
\maketitle
\section{Introduction}
Wagiman (ISO 639-3 waq) is an endangered non-Pama-Nyungan language of the Pine Creek area, Katherine region, Northern Territory. Speaker estimates vary from about ten (mostly elderly) to two.
\section{Primary Educational Resources}
No school primers were found. The online dictionary below is the main teaching-adjacent resource.
\section{Linguistic and Reference Materials}
\begin{itemize}
\item The Wagiman online dictionary (AIATSIS; Stephen Wilson and Mark Harvey with language consultants): \url{https://aiatsis.gov.au/sites/default/files/e_access/digital/a339234/home.html}
\item About the dictionary: \url{https://aiatsis.gov.au/sites/default/files/e_access/digital/a339234/aboutdict.html}
\item The language and its speakers: \url{https://aiatsis.gov.au/sites/default/files/e_access/digital/a339234/langspeakers.html}
\item Austlang entry N27: \url{https://collection.aiatsis.gov.au/austlang/language/N27}
\item Wagiman language, overview: \url{https://en.wikipedia.org/wiki/Wagiman_language}
\end{itemize}
\section{Religious and Community Resources}
Links and references page of the dictionary site: \url{https://aiatsis.gov.au/sites/default/files/e_access/digital/a339234/linksrefs.html}. The site notes community-produced \emph{Wagiman Plants and Animals} and a plants and animals calendar; I found no URLs for these.
\section{Digital and Audio Resources}
\begin{itemize}
\item Three Wagiman texts with English translations (from \emph{Wakgala mahan matjjin Wagiman}, 1998): \url{https://aiatsis.gov.au/sites/default/files/e_access/digital/a339234/texts/texts.html}
\item PARADISEC blog post on delivery of the electronic dictionary: \url{https://www.paradisec.org.au/blog/?p=3840}
\end{itemize}
\section{Curriculum Development Context}
Wagiman is not known to be taught in schools. Documentation (dictionary, texts, Wilson's 1999 \emph{Coverbs and Complex Predicates in Wagiman}) exists, but no curriculum, and any work would need Wagiman community leadership.
\end{document}